\section{An all-degree obstruction for two alternating two-box factors}
\label{sec:alternating-two-factor}

\paragraph{Status and exact scope.}
This is a separate new author derivation, not a change to the frozen
D-ATT-01--03 section. The root and independent /root/imprimitive whole-proof, source and
code readings have passed, with distinct exact supplements.
Keep either actual rational module $W$ in (ATT.1), with $n\ge15$,
and put $a=\dim W$, $V=W\otimes W$, $e=a^2$.
The source theorem below assumes that the inducing image acts
transitively on the two simple ordinary factors $A_n\times A_n$.
It does not assert this transitivity for an arbitrary tensor product.

\begin{lemma}[D-ATT-04: a common index bound]
\label{D-ATT-04}
Put $A=(n-5)/2$ and $B=n/4$. For either module and every
nonidentity permutation $\pi\in S_n$,
\[
                 \nu(W(\pi))\ge A\operatorname{ind}(\pi)+B.
                                                               \tag{ATF.1}
\]
\end{lemma}

\begin{proof}
For UP, the target-one inequality (UP.3) is exactly (ATF.1).
For a nonidentity target scalar, (UP.4) has constant term
$(n-4)/2\ge n/4$, so it is at least as strong.
For EX, the target-one inequality (EX.4) is also stronger,
by $\operatorname{ind}(\pi)/2+n/12$.
Write $k=\operatorname{ind}(\pi)$ for the remaining EX targets.
If $k\ge n/2-2$, (EX.2) exceeds $Ak+B$ by
$k/2+1-n/4\ge0$. If $k<n/2-2$, (EX.3) gives the lower bound
\[
 a-\frac{(n-1)^2}{4}=rac{n^2-4n+3}{4}
 >\frac{n^2-8n+20}{4}
   =A(n/2-2)+B>Ak+B.                                      \tag{ATF.2}
\]
The strict middle inequality holds since $4n-17>0$.
Minimizing over targets proves the assertion. The identity
permutation is excluded from (ATF.1), as it must be.
\end{proof}

\begin{lemma}[D-ATT-05: actual weighted swaps]
\label{D-ATT-05}
Let $S(u\otimes v)=v\otimes u$. For $x,y\in S_n$, the
operator $T=(W(x)\otimes W(y))S$ is conjugate to
$(W(h)\otimes I)S$, where $h=xy$. For every nonzero scalar $\lambda$,
\[
 \operatorname{rank}(1-\lambda T)
       \ge\frac{a^2+a\nu(W(h))-a}{2}.                     \tag{ATF.3}
\]
In particular every scalar times a swap has support at least
$(a-1)e/(2a)$. If $h=1$ and $\lambda\notin\mu_6$, its
actual moving rank is $e$.
\end{lemma}

\begin{proof}
Conjugate by $I\otimes W(y)$. Let $\alpha_1,\ldots,\alpha_a$
be the eigenvalues of $W(h)$, with repetition. In an eigenbasis,
the tensor word $(i,i)$ has eigenvalue $\alpha_i$.
The two words $(i,j),(j,i)$, with $i\ne j$, form a weighted
two-cycle with the two distinct roots of
$z^2=\alpha_i\alpha_j$ as eigenvalues.

Fix a proposed eigenvalue $\zeta$. Let $N_\zeta$ count ordered
pairs with $\alpha_i\alpha_j=\zeta^2$, and let $d_+,d_-$
count $\alpha_i=\zeta,-\zeta$, respectively. Its exact
multiplicity is
\[
             m_\zeta=\frac{N_\zeta+d_+-d_-}{2}
                      \le\frac{N_\zeta+a}{2}.              \tag{ATF.4}
\]
For each $i$ the possible second indices lie in one eigenspace
of $W(h)$, so
$N_\zeta\le a(a-\nu(W(h)))$.
This proves (ATF.3), including off-spectrum targets, by taking
$\zeta=\lambda^{-1}$. Omitting the nonnegative support term
gives the asserted uniform swap bound.
If $h=1$, the conjugated operator is $S$, with eigenvalues
only $1,-1$. A scalar outside $\mu_6$ is neither of these,
so its actual fixed space is zero.
\end{proof}

\begin{theorem}[D-ATT-06: no transitive two-factor source in any degree]
\label{D-ATT-06}
Suppose the faithful finite inducing image $J$ contains the
ordinary irreducibly acting core $A_n\times A_n$ on $W\otimes W$,
where $n\ge15$ and $W$ is either module in (ATT.1), and
conjugation by $J$ is transitive on these two factors.
If the original imaginary quadratic field is
$K\ne\mathbf Q(\sqrt{-3})$, contained in the actual inducing
character field, its six-dimensional multiplicity variation
cannot have period-image rank nine. All actual scalar phases
and both compatible outer cosets on each factor are included.
\end{theorem}

\begin{proof}
The actual intermediate-cover argument in D-ATT-03 again gives
genus zero, a generating product-one tuple of $r\ge12$ genuine
inertias in $J$, and total moving rank $D=2e+6$.
Use D-ATT-02 to write each inertia uniquely as $\lambda T$.
Its permutation of the two factors is either the identity
or a swap. Write $P$ for the number of swaps, and put
\[
 d_0=\begin{cases}n-3,&\mathrm{UP},\\n-2,&\mathrm{EX},\end{cases}
 \qquad \mu=d_0/a,\quad c=b/a,\quad s=(a-1)/(2a).          \tag{ATF.5}
\]
Here $b$ is exactly as in (ATT.1).

Transitivity forces $P>0$, and the product-one relation
forces $P$ even. If $P\ge4$, four swap bounds and eight
other genuine marks give
\[
    D\ge(4s+8\mu)e
      =2e+a(8d_0-2)>2e+6.                                \tag{ATF.6}
\]
Thus $P=2$.

At least two scalar phases lie outside $\mu_6$, by D-ATT-02.
If any such mark were not one of these two swaps, choose
two marked entries and a swap distinct from them. The two
marked bounds, one swap bound, and nine other genuine marks
would imply
\[
 D\ge\left(2\frac{3-c}{4}+s+9\mu\right)e
      =2e+a\left(9d_0-\frac{b+1}{2}\right)>2e+6.           \tag{ATF.7}
\]
The last inequality follows from $a\ge90$, $d_0\ge12$ and
$b\le d_0$, and is valid for both parities. Therefore the
two swaps are precisely the scalar marks outside $\mu_6$.
In particular both satisfy that scalar condition.

For a non-swap inertia write
$g=\lambda(W(x)\otimes W(y))$ and
$I_g=\operatorname{ind}(x)+\operatorname{ind}(y)$.
If at least one of $x,y$ is nonidentity, convolution gives
\[
 \operatorname{rank}(1-g)
  \ge a\max\{\nu(W(x)),\nu(W(y))\}
  \ge\frac a2(AI_g+B),                                  \tag{ATF.8}
\]
since at least one factor supplies the positive constant $B$
in (ATF.1). If both are the identity, genuine inertia means
$g$ is a nonidentity scalar, so its moving rank is $e$,
which is still at least the right side of (ATF.8).

For either swap let $h=xy$ be its cycle product. If $h\ne1$,
(ATF.1) and (ATF.3) imply
\[
 \operatorname{rank}(1-g)
   \ge\frac e2+\frac a2\big(A\operatorname{ind}(h)+B\big)
                         -\frac a2.                       \tag{ATF.9}
\]
If $h=1$, the marked scalar and D-ATT-05 give moving rank
$e$. This also implies (ATF.9), because $B-1\le a$.
Thus each swap really supplies its $B$ term, even with
trivial cycle product; this uses its actual marked phase.

We now use a quotient of the actual cover. The tensor-wreath
action is faithful: its kernel would centralize both simple
factors, and the scalar-intersection argument of D-ATT-02
then makes every constituent permutation trivial. Hence
discarding the unique scalar gives an actual homomorphism
$J\longrightarrow S_n\wr S_2$. Its action on two blocks
of $n$ letters is transitive because $J$ contains $A_n^2$
and has a swap. Take the quotient of the given connected
$J$-cover by the stabilizer of one of these $2n$ letters.
This is a connected degree-$2n$ cover of the sphere.

The index of a non-swap on these letters is exactly $I_g$.
A swap with cycle product $h$ has one cycle of length $2l$
for each length-$l$ cycle of $h$, so its index is
$n+\operatorname{ind}(h)$. Riemann--Hurwitz on this actual
quotient therefore gives
\[
 E:=\sum_{\mathrm{non\!\!-\!swap}} I_g
          +\operatorname{ind}(h_1)+\operatorname{ind}(h_2)
                          \ge2n-2.                        \tag{ATF.10}
\]
No abstract generating tuple is substituted for the actual
inertia, and scalar entries simply have zero permutation index.

Summing (ATF.8)--(ATF.9) over all $r$ entries now yields
\[
 \begin{split}
 D&\ge e+\frac a2(AE+rB)-a\\
  &\ge e+\frac a2\big((n-5)(n-1)+3n\big)-a.
 \end{split}                                               \tag{ATF.11}
\]
The inner polynomial is $n^2-3n+5$. For UP it equals
$2a+5$, while for EX it equals $2a+3$. Consequently
\[
 D\ge\begin{cases}2e+3a/2,&\mathrm{UP},\\
                   2e+a/2,&\mathrm{EX},\end{cases}
                            \qquad D>2e+6.                 \tag{ATF.12}
\]
This contradicts the actual multiplicity-six Euler identity.
\end{proof}

\paragraph{Remaining scope and strict movement.}
The all-degree conclusion requires transitivity on the two
simple factors. It is the actual transitive two-factor case
of the main tensor-core gate, not a general assertion about
all factor-preserving tensor products. It removes, for example,
both labels at $n=100$, $t=2$, over $\mathbf Q(i)$, which
the prior finite ranges and the ALT07 necessary inequality
leave open. It asserts no prior existence of those sources.
Three or more factors are not treated by this addendum:
full tensor cycles of larger length may have a marked scalar
eigenvalue even when their cycle product is the identity,
so the decisive identity-product step cannot be silently
carried over. The original/inducing field distinction and
all earlier conditional inputs are unchanged.
